
Mais tout polynôme @@M0@@ est continu dans la topologie induite par @@M1@@. Donc pour @@M2@@ assez grand on a
@@M3@@
Mais par définition @@M4@@ est un entier pour tout @@M5@@. Il s'ensuit que @@M6@@ pour @@M7@@ assez grand. Donc le polynôme @@M8@@ est constant et on trouve en particulier @@M9@@ i.e. @@M10@@.

Les deux cas a) et b) ci-dessus montrent que la restriction de @@M11@@ à @@M12@@ est constante. On va montrer maintenant que @@M13@@ est constante sur chaque classe de @@M14@@ suivant @@M15@@. Pour tout @@M16@@ on désigne par @@M17@@ l'application définie par
@@M18@@
L'application @@M19@@ est une application polynôme de @@M20@@ dans @@M21@@ qui se factorise par @@M22@@, donc constante d'après ce qu'on a vu ci-dessus, d'où l'assertion que @@M23@@ soit constante sur les classes de @@M24@@ suivant @@M25@@. Il s'ensuit qu'on peut écrire @@M26@@, où @@M27@@ est une application polynôme qui se factorise par @@M28@@. Mais par hypothèse @@M29@@ est un groupe de torsion de sorte que @@M30@@ est constante. Donc @@M31@@ est constante, C.Q.F.D.

\textbf{Lemme 2.1.4.} \uline{Soit} @@M32@@ \uline{un anneau commutatif unitaire et} @@M33@@ \uline{un idéal nilpotent de} @@M34@@\uline{. Soit} @@M35@@ \uline{le sous-groupe multiplicatif de} @@M36@@ \uline{formé des éléments inversibles de} @@M37@@ \uline{de la forme} @@M38@@ \uline{avec} @@M39@@\uline{, et soit} @@M40@@ \uline{l'injection naturelle. Alors l'application composée}
